\section{A four-effect coherence bound}
\label{sec:coherence}

The first step in the operational argument is an inequality for four
arbitrary positive local matrices.  It compares the coherence between
two equal-bit branches with the weights of branches in which the two
bits differ.  The matrices need not form measurements: later they will
be densities at a fixed public transcript, where pointwise
normalization is not required.

We use the represented class $\mathscr C$ of
Definition~\ref{def:represented-class}.  Thus, if $\omega\in\mathscr C$
acts on $\C^p\otimes\C^q$, then for some positive integers $r,s$ it has
a representation
\[
 \omega=(\mathcal L\otimes\mathcal M)(\chi\chi^*),
 \qquad
 \chi\in\C^r\otimes\C^s,
\]
where $\mathcal L:\Mat{r}\to\Mat{p}$ and
$\mathcal M:\Mat{s}\to\Mat{q}$ satisfy that
$\mathcal L,\mathcal L\circ T_r,\mathcal M,\mathcal M\circ T_s$
are completely positive.  Only $\omega$ has trace one; the vector and
the rectangular maps have no normalization requirement.  We will use
that freedom to change coordinates on the input of $\mathcal L$.

\begin{lemma}[Positive block estimate]\label{lem:positive-block}
Let $H$ and $J$ be square matrices, possibly of different sizes, and
let $Z$ have the corresponding rectangular size.  If
\[
 \begin{pmatrix}H&Z\\ Z^*&J\end{pmatrix}\succeq0,
\]
then
\[
 \norm{Z}_1\le\sqrt{\tr H\,\tr J}.
\]
This estimate includes singular $H$ and $J$.
\end{lemma}

\begin{proof}
A spectral decomposition of the positive block matrix expresses it as
a finite sum of outer products of block columns $\binom{a_k}{b_k}$.
Consequently
$H=\sum_k a_ka_k^*$, $J=\sum_k b_kb_k^*$, and
$Z=\sum_k a_kb_k^*$.  Since a rank-one matrix $ab^*$ has trace norm
$\norm{a}\,\norm{b}$, the triangle inequality and Cauchy--Schwarz give
\[
 \norm{Z}_1
 \le\sum_k\norm{a_k}\,\norm{b_k}
 \le\left(\sum_k\norm{a_k}^2\right)^{1/2}
      \left(\sum_k\norm{b_k}^2\right)^{1/2}
 =\sqrt{\tr H\,\tr J}.
\]
No inverse of a diagonal block is used.
\end{proof}

\begin{lemma}[Four-effect inequality]\label{lem:coherence}
Let $\omega\in\mathscr C$ act on $\C^p\otimes\C^q$.  For arbitrary
$X_0,X_1\succeq0$ in $\Mat{p}$ and $Y_0,Y_1\succeq0$ in $\Mat{q}$, put
\[
 D_{ij}=X_i\otimes Y_j,
 \qquad p_{ij}=\tr(D_{ij}\omega),
 \qquad i,j\in\{0,1\}.
\]
Then
\begin{equation}\label{eq:coherence}
 \norm{\sqrt{D_{00}}\,\omega\sqrt{D_{11}}}_1
 \le \sqrt{p_{00}p_{01}}+\sqrt{p_{10}p_{11}}
       +2\sqrt{p_{01}p_{10}}.
\end{equation}
There is no normalization assumption on $X_0,X_1,Y_0,Y_1$.
\end{lemma}

\begin{proof}
Fix a representation of $\omega$ as above.  Pull the four matrices back
to the two input spaces by the Hilbert--Schmidt adjoints:
\[
 A_i=\mathcal L^\dagger(X_i)\in\Mat{r},
 \qquad B_j=\mathcal M^\dagger(Y_j)\in\Mat{s}.
\]
These matrices are positive.  For instance, for every $a\in\C^r$,
\[
 a^*A_i a=\tr\bigl[X_i\mathcal L(aa^*)\bigr]\ge0.
\]
We first replace the representation by one in which $A_0$ and $A_1$ commute;
the two complete-positivity conditions will then give complementary
estimates for each off-diagonal term of the input.

\paragraph{Normalization on the support.}
Set $H=A_0+A_1$, and let $\Pi$ be the orthogonal projection onto
$\ran H$.  If $a\in\ker H$, the nonnegative numbers
$a^*A_i a=\norm{\sqrt{A_i}a}^2$ sum to zero, so each $A_i$ annihilates
$\ker H$.  It follows also that $A_i=\Pi A_i\Pi$.
Define $S$ to be $H^{-1/2}$ on $\ran H$ and the identity on $\ker H$.
The matrix $S$ is positive and invertible on all of $\C^r$, including
when $H=0$.  The two positive matrices
\[
 A_i'=S^*A_iS
\]
are supported on $\Pi$ and satisfy $A_0'+A_1'=\Pi$.  Hence they commute:
$[A_0',A_1']=[A_0',\Pi-A_0']=0$.

This congruence can be made inside the representation without changing
$\omega$.  With $\Ad_S(U)=SUS^*$, replace
\[
 \mathcal L\quad\hbox{by}\quad\mathcal L'=
 \mathcal L\circ\Ad_S,
 \qquad
 \chi\quad\hbox{by}\quad\chi'=(S^{-1}\otimes I_s)\chi.
\]
The resulting output is the same, and the pullbacks of $X_i$ are $A_i'$.
The new map is CP, and its input-transpose condition follows from
\[
 \Ad_S\circ T_r=T_r\circ\Ad_{\overline S},
 \qquad
 \mathcal L'\circ T_r=(\mathcal L\circ T_r)
                         \circ\Ad_{\overline S}.
\]
Thus it remains a representation of the kind required by
Definition~\ref{def:represented-class}.  We now suppress the primes.

\paragraph{The two block estimates.}
Let $C=\mathcal L\otimes\mathcal M$.  For input vectors
$u,z\in\C^r\otimes\C^s$, define
\[
 \begin{split}
 p_{ij}(u)&=u^*(A_i\otimes B_j)u
       =\tr\bigl[D_{ij}C(uu^*)\bigr],\\
 Z(u,z)&=\sqrt{D_{00}}\,C(uz^*)\sqrt{D_{11}}.
 \end{split}
\]
Complete positivity of $C$ means that applying it entrywise to
\[
 \begin{pmatrix}uu^*&uz^*\\ zu^*&zz^*\end{pmatrix}
 =\binom{u}{z}\binom{u}{z}^*
\]
gives a positive block matrix.  Conjugate that matrix by
$\operatorname{diag}(\sqrt{D_{00}},\sqrt{D_{11}})$ and apply
Lemma~\ref{lem:positive-block}.  The result is
\begin{equation}\label{eq:direct-block-bound}
 \norm{Z(u,z)}_1\le\sqrt{p_{00}(u)p_{11}(z)}.
\end{equation}

The second estimate concerns the same $Z(u,z)$, but exchanges the two
input vectors in the diagonal weights.  In the product input basis the
\emph{global input transpose} is $T_{rs}=T_r\otimes T_s$.  The map
\[
 C\circ T_{rs}=(\mathcal L\circ T_r)\otimes
                (\mathcal M\circ T_s)
\]
is CP.  Apply the preceding block argument to this map, with the
ordered pair $(\overline z,\overline u)$ and the same two output filters.
The identities
\[
 (\overline z\,\overline u^*)^\T=uz^*,\qquad
 (\overline z\,\overline z^*)^\T=zz^*,\qquad
 (\overline u\,\overline u^*)^\T=uu^*
\]
show that the off-diagonal output is unchanged while the two diagonal
inputs are exchanged.  We obtain
\begin{equation}\label{eq:reversed-block-bound}
 \norm{Z(u,z)}_1\le\sqrt{p_{00}(z)p_{11}(u)}.
\end{equation}
It is the transpose on both input factors here that supplies this
reverse estimate.

\paragraph{Splitting the input vector.}
Choose a common orthonormal eigenbasis of the commuting $A_0,A_1$.
Let $P$ be the projection onto the basis vectors on which the
eigenvalue of $A_0$ is at least the eigenvalue of $A_1$, and set
\[
 v=(P\otimes I_s)\chi,\qquad
 w=((I_r-P)\otimes I_s)\chi.
\]
Then $\chi=v+w$.  Because $PA_i(I_r-P)=0$, the cross terms in the
quadratic forms defining $p_{ij}$ vanish.  The eigenvalue comparison
on the two subspaces, together with $B_0,B_1\succeq0$, therefore gives
\begin{equation}\label{eq:spectral-weights}
 \begin{gathered}
 p_{ij}=p_{ij}(v)+p_{ij}(w),\\
 p_{11}(v)\le p_{01}(v),\qquad
 p_{00}(w)\le p_{10}(w).
 \end{gathered}
\end{equation}
For example, $A_0-A_1$ is positive on $\ran P$, so its tensor
product with $B_1$ has nonnegative quadratic form on $v$; the second
comparison follows on $\ker P$ using $B_0$.

We can now bound the four terms in the expansion of $C(\chi\chi^*)$.
For the diagonal terms, \eqref{eq:direct-block-bound} and
\eqref{eq:spectral-weights} give
\[
 \begin{aligned}
 \norm{Z(v,v)}_1
 &\le\sqrt{p_{00}(v)p_{11}(v)}
  \le\sqrt{p_{00}p_{01}},\\
 \norm{Z(w,w)}_1
 &\le\sqrt{p_{00}(w)p_{11}(w)}
  \le\sqrt{p_{10}p_{11}}.
 \end{aligned}
\]
For $Z(v,w)$ use the reverse estimate
\eqref{eq:reversed-block-bound}; for $Z(w,v)$ use the direct estimate
\eqref{eq:direct-block-bound}.  Both right sides become the same:
\[
 \begin{aligned}
 \norm{Z(v,w)}_1&\le\sqrt{p_{00}(w)p_{11}(v)}
                       \le\sqrt{p_{10}p_{01}},\\
 \norm{Z(w,v)}_1&\le\sqrt{p_{00}(w)p_{11}(v)}
                       \le\sqrt{p_{10}p_{01}}.
 \end{aligned}
\]
Finally,
\[
 \sqrt{D_{00}}\,\omega\sqrt{D_{11}}
 =Z(v,v)+Z(v,w)+Z(w,v)+Z(w,w).
\]
The triangle inequality now yields \eqref{eq:coherence}.
\end{proof}

By Lemma~\ref{lem:tensor-closure}, the same inequality applies to
arbitrary joint positive matrices on all of Alice's and all of Bob's
copies of a state in $\mathscr C$.  Its lack of any normalization
condition is the feature that will let us apply it at each complete
public transcript.
